\medskip

\textbf{Partie A}

\medskip

On considère une fonction $f$ définie sur $[0~;~ +\infty[$, représentée par la courbe
$\mathcal{C}$ ci-dessous.

La droite $T$ est tangente à la courbe $\mathcal{C}$ au point A d'abscisse $\dfrac52$.

\begin{center}
\psset{xunit=2cm,arrowsize=2pt 3,comma}
\begin{pspicture}(-0.5,-5.5)(5.5,5.5)
\multido{\n=0.0+0.5}{12}{\psline[linewidth=0.2pt](\n,-5.5)(\n,5.5)}
\multido{\n=0.0+0.1}{56}{\psline[linewidth=0.1pt](\n,-5.5)(\n,5.5)}
\multido{\n=-5+1}{11}{\psline[linewidth=0.2pt](0,\n)(5.5,\n)}
\multido{\n=-5.4+0.2}{55}{\psline[linewidth=0.1pt](0,\n)(5.5,\n)}
\psaxes[linewidth=1.25pt,Dx=0.5]{->}(0,0)(0,-5.5)(5.5,5.5)
\psplot[plotpoints=2000,linewidth=1.25pt,linecolor=cyan]{0}{5.5}{4 x mul 2 sub 2.71828 x 1 sub exp div}
\psplotTangent{2.5}{3}{4 x mul 2 sub 2.71828 x 1 sub exp div}
\uput[r](0.25,-3){\cyan $\mathcal{C}$}\uput[u](0.25,3.8){$T$}
\psdots[dotstyle=+,dotscale=2.2](2.5,1.785)\uput[ur](2.5,1.785){A}
\end{pspicture}
\end{center}

\begin{enumerate}
\item Dresser, par lecture graphique, le tableau des variations de la fonction $f$ sur l'intervalle [0~;~5].
\item Que semble présenter la courbe $\mathcal{C}$ au point A ?
\item La dérivée $f'$ et la dérivée seconde $f''$ de la fonction $f$ sont représentées par les courbes ci-dessous.

Associer à chacune de ces deux fonctions la courbe qui la représente.

Ce choix sera justifié.

\begin{tabularx}{\linewidth}{*{2}{>{\centering \arraybackslash}X}}
\psset{xunit=1cm,yunit=0.2cm,arrowsize=2pt 3,comma}
\begin{pspicture}(-0.5,-26)(5.5,0.5)
\multido{\n=0.0+0.5}{12}{\psline[linewidth=0.2pt](\n,-26)(\n,0.5)}
\multido{\n=-26+2}{14}{\psline[linewidth=0.2pt](0,\n)(5.5,\n)}
\psaxes[linewidth=1.25pt,Dx=0.5,Dy=2,labelFontSize=\scriptstyle]{->}(0,0)(-0.5,-26)(5.5,0.5)
\psplot[plotpoints=2000,linewidth=1.25pt,linecolor=blue]{0}{5.5}{4 x mul 10 sub 2.71828 x 1 sub exp div}
\uput[r](0.25,-20){\blue $\mathcal{C}_1$}
\end{pspicture}&
\psset{xunit=1cm,yunit=0.28cm,arrowsize=2pt 3,comma}
\begin{pspicture}(-0.5,-2)(5.5,17)
\multido{\n=0.0+0.5}{12}{\psline[linewidth=0.2pt](\n,-2)(\n,16)}
\multido{\n=-2+2}{10}{\psline[linewidth=0.2pt](0,\n)(5.5,\n)}
\psaxes[linewidth=1.25pt,Dx=0.5,Dy=2,labelFontSize=\scriptstyle]{->}(0,0)(-0.5,-2)(5.5,17)
\psplot[plotpoints=2000,linewidth=1.25pt,linecolor=red]{0}{5.5}{6 4 x mul sub 2.71828 x 1 sub exp div}
\uput[r](0.5,7){\red $\mathcal{C}_2$}
\end{pspicture}\\
Courbe $\mathcal{C}_1$	&Courbe $\mathcal{C}_2$\\
\end{tabularx}

\begin{minipage}{0.48\linewidth}
\item  La courbe $\mathcal{C}_3$ ci-contre peut-elle être la
représentation graphique sur $[0~;~+\infty[$ d'une primitive de la fonction $f$ ? Justifier.
\end{minipage}\hfill
\begin{minipage}{0.48\linewidth}
\psset{xunit=1cm,yunit=0.5cm,arrowsize=2pt 3,comma}
\begin{pspicture}(-0.5,-0.4)(5.5,7.5)
\multido{\n=0.0+0.5}{12}{\psline[linewidth=0.2pt](\n,-0.5)(\n,7.5)}
\multido{\n=0+1}{8}{\psline[linewidth=0.2pt](0,\n)(5.5,\n)}
\psaxes[linewidth=1.25pt,Dx=0.5,Dy=2,labelFontSize=\scriptstyle]{->}(0,0)(0,-0.5)(5.5,7.5)
\psplot[plotpoints=2000,linewidth=1.25pt,linecolor=cyan]{0}{5.5}{4 x mul 2 add 2.71828 x 1 sub exp div }
\uput[ur](1.5,5){\cyan $\mathcal{C}_3$}
\end{pspicture}
\end{minipage}
\end{enumerate}

\medskip

\textbf{Partie B}

\medskip

Dans cette partie, on considère que la fonction $f$, définie et deux fois dérivable sur
$[0~; ~+\infty[$,
est définie par
\[f(x)  = (4x - 2)\e^{-x+1}.\]

On notera respectivement $f'$ et $f''$ la dérivée et la dérivée seconde de la fonction $f$.

\medskip

\begin{enumerate}
\item Étude de la fonction $f$
	\begin{enumerate}
		\item Montrer que $f'(x) = (-4x + 6)\e^{-x+1}$.
		\item Utiliser ce résultat pour déterminer le tableau complet des variations de la fonction $f$ sur $[0~; ~+\infty[$. On admet que $\displaystyle\lim_{x \to + \infty}f(x) = 0$.
		\item Étudier la convexité de la fonction $f$  et préciser l'abscisse d'un éventuel point d'inflexion de la courbe représentative de $f$.
	\end{enumerate}
\item On considère une fonction $F$ définie sur $[0~; ~+\infty[$ par $F(x) = (ax + b)\e^{-x+1}$, où $a$ et $b$ sont deux nombres réels.
	\begin{enumerate}
		\item Déterminer les valeurs des réels $a$ et $b$ telles que la fonction $F$ soit une primitive de la fonction $f$ sur $[0~; ~+\infty[$.
		\item On admet que $F(x) = (-4x - 2)\e^{-x+1}$ est une primitive de la fonction $f$ sur
	$[0~; ~+\infty[$.

En déduire la valeur exacte, puis une valeur approchée à $10^{-2}$ près, de l'intégrale

\[I = \displaystyle\int_{\frac32}^8 f(x)\:\text{d}x.\]
	\end{enumerate}
\begin{minipage}{0.4\linewidth}
\item  Une municipalité a décidé de construire une piste de trottinette freestyle.

Le profil de cette piste est donné par la courbe représentative de la fonction $f$
sur l'intervalle $\left[\frac32~;~ 8\right]$.

L'unité de longueur est le mètre.
\end{minipage}
\begin{minipage}{0.48\linewidth}
\psset{xunit=1.1cm,yunit=1.5cm,arrowsize=2pt 3,comma}
\begin{pspicture}(1.5,-0.5)(8,2.45)
\psaxes[linewidth=1.25pt,Dx=10,Dy=3,labelFontSize=\scriptstyle](1.5,0)(1.5,0)(8,2.45)
\psplot[plotpoints=2000,linewidth=1.25pt]{1.5}{8}{4 x mul 2 sub 2.71828 x 1 sub exp div}
\multido{\n=1.5+0.5,\na=1.5+0.5}{14}{\uput[d](\n,0){\footnotesize \np{\n}}}
\uput[ur](1.5,2.45){D}
\end{pspicture}
\end{minipage}

	\begin{enumerate}
		\item  Donner une valeur approchée au cm près de la hauteur du point de départ D.
		\item  La municipalité a organisé un concours de graffiti pour orner le mur de profil de la piste. L'artiste retenue prévoit de couvrir environ 75\,\% de la surface du mur.

Sachant qu'une bombe aérosol de $150$~mL permet de couvrir une surface de $0,8$ m$^2$, déterminer le nombre de bombes qu'elle devra utiliser pour réaliser cette œuvre.
	\end{enumerate}
\end{enumerate}

\bigskip
